% !TeX root = ../../Book4.tex
% 纲要标题：### 15.4 正合偶与低阶计算
% 纲要标签：b4:sec:1404
% 纲要位置：计划纲要.md，第 3605 行。

\section{正合偶与低阶计算}
\label{b4:sec:1404}

前两节已从滤过复形直接构造谱序列并证明有限收敛。本节用正合偶重新组织同一过程，再考察非零行列很少时的微分、边缘同态和扩张问题。

% 纲要内容：
%
% * 正合偶的构造
% * 两列、三列和第一象限例子
% * 退化与边缘同态
% * 不忽略隐藏延拓问题
%

\subsection{正合偶的构造}
\label{b4:subsec:1404:01}

谱序列还有一种以长正合列为起点的构造。它把“微分后再看商”的过程压缩成三个相互正合的箭头。

\begin{definition}\label{b4:def:exact-couple}
在交换范畴中，给定对象 \(D,E\) 与态射 \(i:D\to D\)、\(j:D\to E\)、\(k:E\to D\)。若
\[
 \operatorname{im}i=\ker j,\quad
 \operatorname{im}j=\ker k,\quad
 \operatorname{im}k=\ker i,
\]
则称这些数据构成\term[zhenghe-ou]{正合偶}[exact couple]。分次情形允许三个态射具有指定的次数。
\end{definition}

三个核像等式要求 \(D\xrightarrow{j}E\xrightarrow{k}D\xrightarrow{i}D\) 在每个相邻位置正合；分次情形中的几个 \(D\) 可以处于不同次数。用于滤过复形时，它把长正合列按这些次数汇总，而不是给出一条普通短正合列。导出一次时，将 \(E\) 换成 \(H(E,jk)\)，将 \(D\) 换成 \(i\) 的像；若新的箭头仍构成正合偶，就能重复这一过程。

\begin{proposition}\label{b4:prop:derived-exact-couple}
设 \(\mathcal A\) 为交换范畴，\((D,E,i,j,k)\) 是 \(\mathcal A\) 中的正合偶，\(i:D\to D\)、\(j:D\to E\)、\(k:E\to D\)。令 \(\mathrm{d}=jk\)，则 \(\mathrm{d}^2=0\)。取
\[
 D'=\operatorname{im}i,\quad E'=\ker \mathrm{d}/\operatorname{im}\mathrm{d},
\quad i'=i|_{D'},\quad j'(ia)=[ja],\quad k'([e])=ke.
\]
这些公式良定义，并给出新的正合偶。
\end{proposition}
\begin{proof}
因为 \(kj=0\)，有 \(\mathrm{d}^2=0\)。\(i\) 保持 \(D'=\operatorname{im}i\)，故有其限制 \(i'\)。由 \(\mathrm{d}j=0\)，\(j\) 先因子化到 \(\ker\mathrm{d}\)，再投影到 \(E'\)。这一复合在 \(\ker i=\operatorname{im}k\) 上为零，因为 \(jk=\mathrm{d}\) 的像在 \(E'\) 中为零；因此它经满态射 \(D\twoheadrightarrow\operatorname{im}i\) 唯一下降为 \(j':D'\to E'\)。另一方面，\(k\) 在 \(\ker\mathrm{d}\) 上的像位于 \(\ker j=D'\)，且它杀掉 \(\operatorname{im}\mathrm{d}\)，因为 \(k\mathrm{d}=kjk=0\)；故也唯一下降为 \(k':E'\to D'\)。这给出了三个实际态射以及所列公式。

三条新箭头的操作不同：\(i'\) 限制定义域，\(k'\) 从余圈代表下降到同调类，\(j'\) 则须先沿 \(i\) 反向取原像，再用 \(j\)，最后取类。用代表记号把最后这条路径写成
\[
ia\ \xleftarrow{\ i\ }\ a
\ \xmapsto{\ j\ }\ ja
\ \xmapsto{\text{取同调类}}\ [ja].
\]
左箭头只表示选择 \(ia\) 的原像，不表示 \(i\) 有逆。\(j'\) 尤其不是 \(j\) 在 \(D'\) 上的限制：由 \(D'=\ker j\)，后者恒为零，而 \(j'\) 记录的是原像 \(a\) 的像。

以下按\cref{b4:lem:generalized-element-calculus}在共同满态射后作有限提升，并使用广义元素记号。若 \(ia=ia'\)，则 \(a-a'=kb\)，所以 \(ja-ja'=\mathrm{d}b\)，说明 \(j'\) 的公式与提升选择无关。若 \(\mathrm{d}e=0\)，则 \(ke\in\ker j=\operatorname{im}i\)；更改 \(e\) 一个 \(\mathrm{d}b\) 不改变 \(ke\)，这也核对了 \(k'\) 的公式。

若 \(j'(ia)=0\)，则 \(ja=jkb\)，于是 \(a-kb=ic\)。对两边施 \(i\)，得到 \(ia-ikb=i^2c\)；正合性给出 \(ik=0\)，故 \(ia=i^2c\)，这证明 \(\ker j'=\operatorname{im}i'\)。若 \(k'[e]=0\)，则 \(ke=0\)，可写 \(e=ja\)，从而 \([e]=j'(ia)\)，给出第二处正合。最后，若 \(ia\in D'\) 且 \(i(ia)=0\)，写 \(ia=ke\)。由于 \(jke=jia=0\)，\(e\) 是 \(\mathrm{d}\)-余圈，且 \(k'[e]=ia\)。各相邻复合为零，有限提升判据因而给出第三处正合。
\end{proof}

对短正合复形列 \(0\to F^{p+1}C\to F^pC\to\operatorname{gr}_F^pC\to0\)，\cref{b4:thm:cohomology-long-exact}给出的相邻四项为
\[
 H^n(F^{p+1}C)\xrightarrow{i}H^n(F^pC)
 \xrightarrow{j}H^n(\operatorname{gr}_F^pC)
 \xrightarrow{k}H^{n+1}(F^{p+1}C).
\]
令 \(n=p+q\)，记
\[
 D_1^{p,q}=H^{p+q}(F^pC),\qquad E_1^{p,q}=H^{p+q}(\operatorname{gr}_F^pC).
\]
三个箭头依次为
\[
 i:D_1^{p+1,q-1}\to D_1^{p,q},\quad
 j:D_1^{p,q}\to E_1^{p,q},\quad
 k:E_1^{p,q}\to D_1^{p+1,q}.
\]
先施连接态射 \(k\)，再施下一层的商所诱导的 \(j\)，便得到 \(jk:E_1^{p,q}\to E_1^{p+1,q}\)。将这条短正合复形列映到\secref{b4:sec:1402}中的相邻滤过商短正合列，左端和中间项取模 \(F^{p+2}C\)，右端保持不变；连接态射的自然性说明 \(jk\) 恰是那里定义的 \(\mathrm d_1\)。

\ProseParagraphBreak
第一次导出时，\(D_2^{p,q}\) 是 \(H^{p+q}(F^{p+1}C)\to H^{p+q}(F^pC)\) 的像。对其中的 \(ia\)，取原像要回到 \(D_1^{p+1,q-1}\)，然后 \(j\) 把它送到 \(E_1^{p+1,q-1}\)，取类便得到 \(j_2(ia)\in E_2^{p+1,q-1}\)。因此这一回到较深层的步骤使 \(j_2\) 的次数成为 \((1,-1)\)；\(k_2\) 仍具有次数 \((1,0)\)，于是 \(j_2k_2=\mathrm{d}_2\) 具有次数 \((2,-1)\)。后续导出反复进行同样的原像选择，下面把它与代表元公式逐项比较。

\begin{proposition}\label{b4:prop:filtered-exact-couple-comparison}
设 \(\mathcal A\) 为交换范畴，\((C,F)\) 为 \(\mathcal A\) 中的递减滤过上链复形。由各短正合列
\(0\to F^{p+1}C\to F^pC\to\operatorname{gr}_F^pC\to0\)
的长正合列组成上述正合偶，并将与第 \(r\) 页相配的正合偶中的对象记为 \(D_r,\widetilde E_r\)，其中 \(r\ge1\)。则
\[
D_r^{p,q}=
\operatorname{im}\bigl(H^{p+q}(F^{p+r-1}C)
\longrightarrow H^{p+q}(F^pC)\bigr),
\]
且存在对滤过复形映射自然的同构
\[
\theta_r:E_r^{p,q}\xrightarrow{\ \sim\ }\widetilde E_r^{p,q},
\]
其中 \(E_r\) 是\cref{b4:eq:spectral-page-formula}的直接构造。它们保持 \(\mathrm{d}_r\)，并将直接构造中指定的
\(H(E_r,\mathrm{d}_r)\simeq E_{r+1}\)
与导出正合偶的页间同调识别对应起来。
\end{proposition}
\begin{proof}
双分次对象按分量族处理，以下暂省略双次数，所有公式均逐分量使用。记第一正合偶的三个箭头为 \(i,j,k\)，并沿用\cref{b4:lem:generalized-element-calculus}的广义元素及共同满态射后有限提升约定。反复导出所得的对象可写成
\[
D_r=\operatorname{im}i^{r-1},\qquad
\widetilde E_r=N_r/B_r,
\quad N_r=k^{-1}(\operatorname{im}i^{r-1}),\quad
B_r=j(\ker i^{r-1}).
\]
这里 \(N_r\) 要求连接像 \(ke\) 能沿 \(i\) 反向提升 \(r-1\) 次；\(B_r\) 则由被 \(i^{r-1}\) 零化的原像经 \(j\) 得到。当 \(r=1\) 时，\(N_1=E_1\)、\(B_1=0\)；当 \(r=2\) 时，由 \(\operatorname{im}i=\ker j\)、\(\ker i=\operatorname{im}k\)，有
\[
 N_2=k^{-1}(\ker j)=\ker(jk),\qquad
 B_2=j(\operatorname{im}k)=\operatorname{im}(jk),
\]
所以 \(N_2/B_2\) 恰是第一页的同调。对一般的 \(e\in N_r\)，选择 \(a\) 使 \(i^{r-1}a=ke\)，则导出箭头满足
\[
i_r=i|_{D_r},\qquad
j_r(i^{r-1}a)=[ja],\qquad
k_r([e])=ke,\qquad
\widetilde{\mathrm{d}}_r[e]=[ja].
\]
不同选择的 \(a\) 相差 \(\ker i^{r-1}\)，故最后一式的差属于 \(B_r\)。

为证明这些公式在下一次导出后仍成立，先计算余圈。\([e]\) 的微分为零，当且仅当存在 \(b\in\ker i^{r-1}\) 使 \(ja=jb\)。由 \(\ker j=\operatorname{im}i\)，可写 \(a-b=ic\)，于是
\(ke=i^{r-1}a=i^rc\)。反过来，若 \(ke=i^rc\)，选择 \(a=ic\) 即得 \(ja=0\)。所以微分的核是 \(N_{r+1}/B_r\)。

再计算余边。若 \(i^{r-1}a=ke\)，则 \(i^ra=ike=0\)，因此微分像 \([ja]\) 属于 \(B_{r+1}/B_r\)。反过来，若 \(a\in\ker i^r\)，则 \(i^{r-1}a\in\ker i=\operatorname{im}k\)，可以选择 \(e\) 使 \(ke=i^{r-1}a\)。此时 \(e\in N_r\)，且其微分就是 \([ja]\)。因此像恰为 \(B_{r+1}/B_r\)，页间同调恰为 \(N_{r+1}/B_{r+1}\)。同时，导出定义把 \(D_r\) 换成 \(\operatorname{im}i^r\)，并给出相同的 \(j_{r+1},k_{r+1}\) 公式。归纳证明了上述全部描述。

恢复双次数后，\(i\) 来自滤过包含，故其反复像给出命题中的 \(D_r^{p,q}\)。三个箭头的次数为
\[
\begin{aligned}
i_r&:D_r^{p,q}\longrightarrow D_r^{p-1,q+1},\\
j_r&:D_r^{p,q}\longrightarrow\widetilde E_r^{p+r-1,q-r+1},\\
k_r&:\widetilde E_r^{p,q}\longrightarrow D_r^{p+1,q}.
\end{aligned}
\]
因此 \(\widetilde{\mathrm{d}}_r=j_rk_r\) 的次数是 \((r,1-r)\)。

令 \(n=p+q\)。现在以代表元的第一页类比较两种构造：\(x\in Z_r^{p,q}\) 的第一页类将落在 \(N_r^{p,q}\)，模去 \(B_r^{p,q}\) 后应当恰好得到直接构造的第 \(r\) 页。在第一页，\(e\in E_1^{p,q}\) 由 \(x\in F^pC^n\) 表示，且 \(\mathrm{d}x\in F^{p+1}C^{n+1}\)；连接同态 \(ke\) 是 \(\mathrm{d}x\) 在 \(H^{n+1}(F^{p+1}C)\) 中的类。条件 \(e\in N_r^{p,q}\) 要求这一类来自 \(H^{n+1}(F^{p+r}C)\)。这等价于可以写
\[
\mathrm{d}x=y+\mathrm{d}v,
\qquad y\in F^{p+r}C^{n+1},\quad
\mathrm{d}y=0,\quad v\in F^{p+1}C^n.
\]
用 \(x-v\) 代替 \(x\) 不改变第一页的类，并使其微分落入 \(F^{p+r}\)。因此 \(Z_r^{p,q}\) 自然满射到 \(N_r^{p,q}\)，进而满射到 \(N_r^{p,q}/B_r^{p,q}\)。

现在确定这个满射的核。若 \(x\in Z_r^{p,q}\) 的第一页类属于 \(B_r^{p,q}\)，则存在余圈 \(z\in F^pC^n\)，使 \([z]\in H^n(F^pC)\) 在包含到 \(H^n(F^{p-r+1}C)\) 后为零，并使 \(x,z\) 在第一页表示同一类。因此可取
\[
z=\mathrm{d}w,\qquad w\in F^{p-r+1}C^{n-1},
\qquad x-z=u+\mathrm{d}v,
\quad u\in F^{p+1}C^n,\quad v\in F^pC^{n-1}.
\]
由 \(\mathrm{d}u=\mathrm{d}x\in F^{p+r}\)，得到
\(u\in Z_{r-1}^{p+1,q-1}\)。又
\(w+v\in F^{p-r+1}C^{n-1}\)，且
\(\mathrm{d}(w+v)=x-u\in F^pC^n\)，所以
\(w+v\in Z_{r-1}^{p-r+1,q+r-2}\)。于是核包含于
\[
Z_{r-1}^{p+1,q-1}
+\mathrm{d}Z_{r-1}^{p-r+1,q+r-2}.
\]
反过来，第一项在第一页已经为零；第二项由较浅层的余边给出，作为 \(H^n(F^pC)\) 的类正位于 \(\ker i^{r-1}\)，故其第一页类属于 \(B_r\)。这证明核恰为所写分母，得到 \(\theta_r\)。

若 \(x\in Z_r^{p,q}\)，则 \(\mathrm{d}x\) 已在 \(F^{p+r}C^{n+1}\) 中闭合，可以直接用它提升 \(k\) 的像；随后施 \(j_r\)，所得正是 \(\mathrm{d}x\) 的第 \(r\) 页类。因此
\(\theta_r\mathrm{d}_r=\widetilde{\mathrm{d}}_r\theta_r\)。当 \(x\in Z_{r+1}\) 时，上面将余圈送到 \(N_{r+1}/B_{r+1}\) 的同调识别，仍以同一个 \(x\) 表示下一页的类；这与直接构造中指定的页间识别相同。滤过复形映射保持所有包含、连接、核、像和商，故全部同构均自然。
\end{proof}

这个比较的同调型写法见 Weibel\cite[Theorem 5.9.4, pp. 155--156]{Weibel1994HomologicalAlgebra}；这里的递减滤过与上同调双次数给出了相应的箭头方向。

\subsection{两列、三列与第一象限例子}
\label{b4:subsec:1404:02}

设第一象限谱序列强收敛于 \(H^*\)。若其第二页只在 \(p=0,1\) 两列非零，则从 \(r=2\) 起微分都不可能连接两个非零列，故 \(E_2=E_\infty\)。对总次数 \(n\)，收敛给出
\[
 0\longrightarrow E_2^{1,n-1}\longrightarrow H^n
 \longrightarrow E_2^{0,n}\longrightarrow0.
\]
这说明两列计算天然产生一个扩张，并不自动产生两项直和。

\ProseParagraphBreak
若只在 \(p=0,1,2\) 三列非零，第一个仍可能出现的高阶微分是
\[
 \mathrm{d}_2:E_2^{0,q}\longrightarrow E_2^{2,q-1};
\]
它之后才有 \(E_3=E_\infty\)。\secref{b4:sec:1405}将把这个箭头真正算出来。若谱序列没有列数上界但在第一象限内，对固定 \((p,q)\) 的箭头仍会终止：当 \(r>q+1\) 时出箭头的靶在负行，当 \(r>p\) 时入箭头的源在负列。因此逐位置稳定，并不意味着整个谱序列在同一个有限页退化。

\subsection{退化与边缘同态}
\label{b4:subsec:1404:03}

设第一象限谱序列强收敛于带滤过 \(F\) 的 \(H^*\)。其边缘位置能与目标的端部相联系，理由是目标滤过逐度有限。固定总次数 \(n\)，第一象限条件给出 \(E_\infty^{p,n-p}=0\) 对 \(p<0\) 或 \(p>n\) 成立，因而这些位置的相邻滤过商都为零。从穷尽的浅端逐层向上，得到 \(F^0H^n=H^n\)；从为零的深端逐层向下，得到 \(F^{n+1}H^n=0\)。因此
\(E_\infty^{n,0}=F^nH^n\)，而
\(E_\infty^{0,n}=H^n/F^1H^n\)。

\begin{definition}\label{b4:def:spectral-degeneration}
设 \(\mathcal A\) 为交换范畴，\(E_r^{p,q}\) 为 \(\mathcal A\) 中的上同调型谱序列，\(s\) 为其一个页数。若从页 \(s\) 开始的所有微分均为零，则称谱序列在第 \(s\) 页\term[tuihua-puxulie]{退化}[degeneration]。若谱序列位于第一象限并强收敛于带滤过 \(F\) 的 \(H^*\)，则对 \(n\ge0\)，目标滤过的端部诱导
\[
 E_2^{n,0}\longrightarrow E_\infty^{n,0}=F^nH^n\longrightarrow H^n,
 \qquad
 H^n\longrightarrow H^n/F^1H^n=E_\infty^{0,n}\longrightarrow E_2^{0,n},
\]
称为相应的\term[bianyuan-tongtai]{边缘同态}[edge homomorphisms]。
\end{definition}

第一条边缘同态由各页的商映射接上 \(F^nH^n\hookrightarrow H^n\) 组成，因为底行从第二页起只有入箭头；第二条由 \(H^n\twoheadrightarrow H^n/F^1H^n\) 接上各页的子对象包含组成，因为最左列只有出箭头。若谱序列从第零页给出，第一、二页仍可有竖直或水平微分，故这个端部描述从第二页开始。

\begin{proposition}\label{b4:prop:spectral-single-row}
设交换范畴中的第一象限上同调型谱序列强收敛于 \(H^*\)。若其第二页只在 \(q=0\) 行非零，则边缘同态 \(E_2^{n,0}\to H^n\) 对每个 \(n\ge0\) 为同构；若只在 \(p=0\) 列非零，则另一边缘同态 \(H^n\to E_2^{0,n}\) 对每个 \(n\ge0\) 为同构。
\end{proposition}
\begin{proof}
在第一种情形，第 \(r\ge2\) 页的出箭头靶行是 \(1-r<0\)，入箭头源在正行却为零，因此全部微分消失。目标滤过只有第 \(n\) 个相邻商可能非零；其余浅层商都为零，故 \(F^nH^n=F^0H^n=H^n\)，而上面的端部识别给出 \(F^{n+1}H^n=0\)。于是边缘同态把 \(E_2^{n,0}=E_\infty^{n,0}\) 同构地映到 \(H^n\)。

在第二种情形，非零项的出箭头落在正列，入箭头来自负列，两者均为零，故也有 \(E_2=E_\infty\)。只有第零个相邻滤过商可能非零，从深端 \(F^{n+1}H^n=0\) 逐层向下得到 \(F^1H^n=0\)。因此另一个边缘同态就是
\(H^n\simeq H^n/F^1H^n=E_\infty^{0,n}=E_2^{0,n}\)。
\end{proof}

\begin{proposition}\label{b4:prop:spectral-five-term}
设交换范畴中的第一象限上同调谱序列从第二页起给出，并强收敛于带有限滤过的 \(H^*\)。则有自然正合列
\[
 0\longrightarrow E_2^{1,0}\longrightarrow H^1
 \longrightarrow E_2^{0,1}\xrightarrow{\mathrm{d}_2}E_2^{2,0}
 \longrightarrow H^2.
\]
与 \(H^1\) 相邻的两条箭头及最右侧箭头是上述边缘同态，最右端不宣称为满态射。
\end{proposition}
\begin{proof}
位置 \((1,0)\) 没有第 \(r\ge2\) 页的入箭头或出箭头，所以 \(E_\infty^{1,0}=E_2^{1,0}\)。位置 \((0,1)\) 只有一条可能的箭头 \(\mathrm{d}_2\) 指向 \((2,0)\)，所以 \(E_\infty^{0,1}=\ker \mathrm{d}_2\)。总次数一的滤过给出 \(0\to E_\infty^{1,0}\to H^1\to E_\infty^{0,1}\to0\)，再包含到 \(E_2^{0,1}\)，得到前半列的正合性。

位置 \((2,0)\) 没有出箭头，唯一可能的入箭头也是这条 \(\mathrm{d}_2\)，故 \(E_\infty^{2,0}=\operatorname{coker}\mathrm{d}_2\)。总次数二的最深可能非零层为 \(F^2H^2\)，并有 \(F^3H^2=0\)，所以 \(E_\infty^{2,0}=F^2H^2\hookrightarrow H^2\)。这给出余下正合性，且最右箭头的像恰为 \(F^2H^2\)。其余相邻商 \(E_\infty^{1,1}\)、\(E_\infty^{0,2}\) 仍可能非零，所以这条箭头不一定覆盖整个 \(H^2\)。全部态射来自滤过及谱序列的自然态射，因此自然。
\end{proof}

\subsection{隐藏延拓问题}
\label{b4:subsec:1404:04}

使用谱序列时可以把计算分成两件事：先求各页及其微分，得到 \(E_\infty\)；再利用目标的其他性质恢复滤过对象。后一步可能使用已知的生成元、乘法、指数或 \(\operatorname{Ext}^1\) 中的扩张类。它不属于再计算一个未出现的 \(\mathrm{d}_r\)。

\ProseParagraphBreak
例如已经知道有限交换群 \(H\) 的两个相邻商均为 \(\mathbb Z/2\mathbb Z\)，则 \(|H|=4\) 已经确定；然而 \(H\) 的指数仍可能是二或四。若此外存在一个四阶元素，则 \(H\cong\mathbb Z/4\mathbb Z\)；若 \(2H=0\)，则 \(H\cong(\mathbb Z/2\mathbb Z)^2\)。判断原对象的指数，需要相邻商以外的扩张信息。
